% GENERATED FILE. Edit canonical lessons, then run npm run generate:books.
% canonical-source-hash: fnv1a64:2f659371d83f4ba0
% canonical-lessons: UR-C329-afsana, UR-C329-musannif, UR-C329-kirdar, UR-C329-bab, UR-C329-anjam

\chapter{A Story in a Book}
\label{ch:ur-a-story-in-a-book}
\hlchaptermodality{329}

\begin{chapteropening}
\textbf{By the end of this chapter:} \emph{I can say what a story is about and who is in it.}

Complete the last lesson of chapter 329: I can say what a story is about and who is in it.
\end{chapteropening}

\section[\texorpdfstring{\ur{افسانہ} (afsānā)}{افسانہ (afsānā)}]{\ur{افسانہ} (afsānā) — a short story}

\label{lesson:UR-C329-afsana}



\begin{quote}
Before the new one: say the Urdu for deceit, then the Urdu for to lead astray.
\end{quote}

\subsection*{You'll want to know: \ur{افسانہ}}
\textbf{\ur{افسانہ}} — \emph{afsānā} — \textquotedblleft{}a short story\textquotedblright{}.

\begin{quote}
\ur{میں} \ur{نے} \ur{کل} \ur{ایک} \ur{افسانہ} \ur{پڑھا۔}
\end{quote}

\begin{quote}
\emph{maiṅ ne kal ek afsānā paṛhā.}
\end{quote}

\begin{quote}
I read a short story yesterday.
\end{quote}

Masculine.

\begin{cousinweb}
\textbf{\ur{افسانہ}} (\emph{afsānā}) is usually traced to Persian \emph{afsāna}, \textquotedblleft{}a tale, a fable\textquotedblright{}; in Urdu writing it is the short story.
\end{cousinweb}

\begin{grammarlens}[title={What you've built}]
One more word for telling what happened.
\end{grammarlens}

\subsection*{Your turn}
\begin{itemize}
\raggedright
  \item \emph{Say it:} \emph{afsānā}
  \item \emph{Say it:} \emph{afsānā}, once more
  \item \emph{Say it:} \emph{maiṅ ne kal ek afsānā paṛhā.} — I read a short story yesterday.
  \item \emph{Recall:} say the Urdu for deceit, then the Urdu for to lead astray, then say \emph{afsānā} again
\end{itemize}

\begin{tcolorbox}[breakable,colback=teal!4,colframe=teal!35!black,title={Before you move on}]
What does \ur{افسانہ} mean? (\textquotedblleft{}A short story\textquotedblright{}.) Say it once more.
\end{tcolorbox}
\section[\texorpdfstring{\ur{مصنف} (musannif)}{مصنف (musannif)}]{\ur{مصنف} (musannif) — an author}

\label{lesson:UR-C329-musannif}



\begin{quote}
Before the new one: say the Urdu for to lead astray, then the Urdu for a short story.
\end{quote}

\subsection*{You'll want to know: \ur{مصنف}}
\textbf{\ur{مصنف}} — \emph{musannif} — \textquotedblleft{}an author\textquotedblright{}.

\begin{quote}
\ur{اس} \ur{افسانے} \ur{کا} \ur{مصنف} \ur{مشہور} \ur{ہے۔}
\end{quote}

\begin{quote}
\emph{is afsāne kā musannif mashhūr hai.}
\end{quote}

\begin{quote}
The author of this story is famous.
\end{quote}

Masculine or feminine, by the person meant.

\begin{cousinweb}
\textbf{\ur{مصنف}} (\emph{musannif}) is usually traced, through Persian, to Arabic \emph{muṣannif}, \textquotedblleft{}one who composes a book\textquotedblright{}.
\end{cousinweb}

\begin{grammarlens}[title={What you've built}]
One more word for telling what happened.
\end{grammarlens}

\subsection*{Your turn}
\begin{itemize}
\raggedright
  \item \emph{Say it:} \emph{musannif}
  \item \emph{Say it:} \emph{musannif}, once more
  \item \emph{Say it:} \emph{is afsāne kā musannif mashhūr hai.} — The author of this story is famous.
  \item \emph{Recall:} say the Urdu for to lead astray, then the Urdu for a short story, then say \emph{musannif} again
\end{itemize}

\begin{tcolorbox}[breakable,colback=teal!4,colframe=teal!35!black,title={Before you move on}]
What does \ur{مصنف} mean? (\textquotedblleft{}An author\textquotedblright{}.) Say it once more.
\end{tcolorbox}
\section[\texorpdfstring{\ur{کردار} (kirdār)}{کردار (kirdār)}]{\ur{کردار} (kirdār) — a character (in a story)}

\label{lesson:UR-C329-kirdar}



\begin{quote}
Before the new one: say the Urdu for a short story, then the Urdu for an author.
\end{quote}

\subsection*{You'll want to know: \ur{کردار}}
\textbf{\ur{کردار}} — \emph{kirdār} — \textquotedblleft{}a character (in a story)\textquotedblright{}.

\begin{quote}
\ur{اس} \ur{ناول} \ur{میں} \ur{دس} \ur{کردار} \ur{ہیں۔}
\end{quote}

\begin{quote}
\emph{is nāvil meṅ das kirdār haiṅ.}
\end{quote}

\begin{quote}
There are ten characters in this novel.
\end{quote}

Masculine.

\begin{cousinweb}
\textbf{\ur{کردار}} (\emph{kirdār}) is usually traced to Persian \emph{kirdār}, \textquotedblleft{}a deed, conduct\textquotedblright{}; the sense of a person in a story is thought to have come from the part one plays.
\end{cousinweb}

\begin{grammarlens}[title={What you've built}]
One more word for telling what happened.
\end{grammarlens}

\subsection*{Your turn}
\begin{itemize}
\raggedright
  \item \emph{Say it:} \emph{kirdār}
  \item \emph{Say it:} \emph{kirdār}, once more
  \item \emph{Say it:} \emph{is nāvil meṅ das kirdār haiṅ.} — There are ten characters in this novel.
  \item \emph{Recall:} say the Urdu for a short story, then the Urdu for an author, then say \emph{kirdār} again
\end{itemize}

\begin{tcolorbox}[breakable,colback=teal!4,colframe=teal!35!black,title={Before you move on}]
What does \ur{کردار} mean? (\textquotedblleft{}A character (in a story)\textquotedblright{}.) Say it once more.
\end{tcolorbox}
\section[\texorpdfstring{\ur{باب} (bāb)}{باب (bāb)}]{\ur{باب} (bāb) — a chapter of a book}

\label{lesson:UR-C329-bab}



\begin{quote}
Before the new one: say the Urdu for an author, then the Urdu for a character (in a story).
\end{quote}

\subsection*{You'll want to know: \ur{باب}}
\textbf{\ur{باب}} — \emph{bāb} — \textquotedblleft{}a chapter of a book\textquotedblright{}.

\begin{quote}
\ur{پہلا} \ur{باب} \ur{سب} \ur{سے} \ur{اچھا} \ur{ہے۔}
\end{quote}

\begin{quote}
\emph{pehlā bāb sab se achchhā hai.}
\end{quote}

\begin{quote}
The first chapter is the best.
\end{quote}

Masculine.

\begin{cousinweb}
\textbf{\ur{باب}} (\emph{bāb}) is usually traced, through Persian, to Arabic \emph{bāb}, \textquotedblleft{}a door\textquotedblright{}, and so a way into one part of a book.
\end{cousinweb}

\begin{grammarlens}[title={What you've built}]
One more word for telling what happened.
\end{grammarlens}

\subsection*{Your turn}
\begin{itemize}
\raggedright
  \item \emph{Say it:} \emph{bāb}
  \item \emph{Say it:} \emph{bāb}, once more
  \item \emph{Say it:} \emph{pehlā bāb sab se achchhā hai.} — The first chapter is the best.
  \item \emph{Recall:} say the Urdu for an author, then the Urdu for a character (in a story), then say \emph{bāb} again
\end{itemize}

\begin{tcolorbox}[breakable,colback=teal!4,colframe=teal!35!black,title={Before you move on}]
What does \ur{باب} mean? (\textquotedblleft{}A chapter of a book\textquotedblright{}.) Say it once more.
\end{tcolorbox}
\section[\texorpdfstring{\ur{انجام} (anjām)}{انجام (anjām)}]{\ur{انجام} (anjām) — an ending; an outcome}

\label{lesson:UR-C329-anjam}



\begin{quote}
Before the new one: say the Urdu for a character (in a story), then the Urdu for a chapter of a book.
\end{quote}

\subsection*{You'll want to know: \ur{انجام}}
\textbf{\ur{انجام}} — \emph{anjām} — \textquotedblleft{}an ending; an outcome\textquotedblright{}.

\begin{quote}
\ur{ناول} \ur{کا} \ur{انجام} \ur{اداس} \ur{تھا۔}
\end{quote}

\begin{quote}
\emph{nāvil kā anjām udās thā.}
\end{quote}

\begin{quote}
The novel had a sad ending.
\end{quote}

Masculine.

\begin{cousinweb}
\textbf{\ur{انجام}} (\emph{anjām}) is usually traced to Persian \emph{anjām}, \textquotedblleft{}an end, a completion\textquotedblright{}. \textbf{\ur{انجام} \ur{کیا} \ur{ہوا؟}} is how did it end?
\end{cousinweb}

\begin{grammarlens}[title={What you've built}]
One more word for telling what happened.
\end{grammarlens}

\subsection*{Your turn}
\begin{itemize}
\raggedright
  \item \emph{Say it:} \emph{anjām}
  \item \emph{Say it:} \emph{anjām}, once more
  \item \emph{Say it:} \emph{nāvil kā anjām udās thā.} — The novel had a sad ending.
  \item \emph{Recall:} say the Urdu for a character (in a story), then the Urdu for a chapter of a book, then say \emph{anjām} again
\end{itemize}

\begin{tcolorbox}[breakable,colback=teal!4,colframe=teal!35!black,title={Before you move on}]
What does \ur{انجام} mean? (\textquotedblleft{}An ending; an outcome\textquotedblright{}.) Say it once more.
\end{tcolorbox}
